\documentclass[12pt,a4paper]{article}

\usepackage{fontspec}
\usepackage{polyglossia}
\setmainlanguage{vietnamese}
\setmainfont{TeX Gyre Termes} % <= Sẵn trong TeX Live, tương đương Times

\usepackage{amsmath,amssymb}
\usepackage{geometry}
\usepackage{fancyhdr}
\usepackage{titlesec}
\usepackage{enumitem}
\usepackage{hyperref}
\usepackage{graphicx}
\usepackage{xcolor}
\usepackage{tabularx}
\usepackage[most]{tcolorbox}
\tcbset{
  colframe=black!40!black,
  colback=gray!2,
  boxrule=0.4pt,
  arc=1pt,
  left=4pt,
  right=4pt,
  top=2pt,
  bottom=2pt,
}

% Watermark settings
% \usepackage{draftwatermark}
% \SetWatermarkText{TMath Coding}     % nội dung watermark
% \SetWatermarkScale{0.75}             % độ lớn (tùy chỉnh)
% \SetWatermarkColor[gray]{0.92}       % màu xám nhạt (0.9 = mờ)
% End watermark settings

\geometry{top=2.5cm, bottom=2.5cm, left=2.5cm, right=2.5cm}
\setlength{\parskip}{6pt}
\setlength{\parindent}{0pt}
\setlength{\headheight}{15pt}
\pagestyle{fancy}
\fancyhf{}
\rhead{Phương trình vi phân}
\lhead{\today}
\rfoot{\thepage}

\titleformat{\section}{\large\bfseries}{\thesection.}{0.5em}{}

\begin{document}

\section{Khái niệm chung}

\subsection{Định nghĩa}
Phương trình vi phân (PTVP) là phương trình có dạng:
\[
F(x, y, y', y'', \dots, y^{(n)}) = 0.
\]
Trong đó $x$ là biến độc lập, $y = y(x)$ là hàm cần tìm.  
Cấp cao nhất của đạo hàm trong phương trình gọi là \textbf{cấp của PTVP}.

\subsection{Nghiệm, nghiệm tổng quát, nghiệm riêng}
\begin{itemize}
    \item \textbf{Nghiệm}: hàm $y(x)$ thỏa phương trình.
    \item \textbf{Nghiệm tổng quát}: 
    \[
    y = \varphi(x, C)
    \]
    chứa hằng số tùy ý.
    \item \textbf{Nghiệm riêng}: nghiệm thu được khi cho $C = C_0$.
\end{itemize}

\subsection{Định lý tồn tại và duy nhất (Cauchy)}
Xét phương trình cấp một:
\[
y' = f(x, y), \qquad y(x_0) = y_0.
\]
Nếu $f$ và $\partial f/\partial y$ liên tục trên miền $D$ thì tồn tại duy nhất nghiệm trong lân cận $x_0$.

\section{Phương trình vi phân cấp một}

\subsection{Phương trình biến số phân li}
\[
f(x)\,dx = g(y)\,dy.
\]
\textbf{Tích phân hai vế:}
\[
\int f(x)\,dx = \int g(y)\,dy + C.
\]

\subsubsection*{Ví dụ}
Giải phương trình:
\[
y' = 2xy.
\]
\textbf{Giải:}
\[
\frac{dy}{y} = 2x\,dx \Rightarrow \ln|y| = x^2 + C.
\]
\[
\Rightarrow y = Ce^{x^2}.
\]

\subsection{Phương trình vi phân đẳng cấp}

\subsubsection*{Định nghĩa}
Phương trình vi phân cấp một
\[
P(x,y)\,dx + Q(x,y)\,dy = 0
\]
được gọi là \textbf{đẳng cấp bậc $n$} nếu cả hai hàm $P(x,y)$ và $Q(x,y)$ là các hàm đẳng cấp bậc $n$, tức là:
\[
P(tx,ty) = t^{n} P(x,y), \qquad 
Q(tx,ty) = t^{n} Q(x,y).
\]

\subsubsection*{Dạng đặc trưng}
Trong phần lớn các bài toán, phương trình đẳng cấp có thể viết dưới dạng:
\[
y' = f\!\left(\frac{y}{x}\right)
\]
hoặc
\[
M(x,y)dx + N(x,y)dy = 0 \quad \text{với $M,N$ đẳng cấp cùng bậc}.
\]

\subsubsection*{Phương pháp đổi biến chuẩn}
Đặt
\[
v = \frac{y}{x} \qquad\Leftrightarrow\qquad y = vx.
\]
Khi đó
\[
y' = v + x v'.
\]

Thay vào phương trình đẳng cấp, ta nhận được phương trình tách biến theo biến $v$ và $x$:
\[
v + xv' = f(v)
\quad\Longrightarrow\quad
x\,\frac{dv}{dx} = f(v) - v.
\]

Đây là phương trình tách biến:
\[
\frac{dv}{f(v)-v} = \frac{dx}{x}.
\]

Tích phân hai vế:
\[
\int \frac{dv}{f(v)-v} = \ln|x| + C.
\]

Trở lại $v = y/x$ ta thu được nghiệm ẩn của phương trình.

\subsubsection*{Ví dụ}

\paragraph{Ví dụ 1.}
Giải phương trình:
\[
\frac{dy}{dx} = \frac{x^2 + y^2}{xy}.
\]

\textbf{Nhận xét:} Vế phải chỉ phụ thuộc vào $y/x$:
\[
\frac{x^2 + y^2}{xy}
= \frac{x^2}{xy} + \frac{y^2}{xy}
= \frac{x}{y} + \frac{y}{x}
= \frac{1}{v} + v,\qquad v=\frac{y}{x}.
\]

Đặt $y = vx$ ⇒ $y' = v + x v'$.

Thay vào:
\[
v + xv' = v + \frac{1}{v}.
\]

Suy ra:
\[
x v' = \frac{1}{v}
\quad\Rightarrow\quad
v\, dv = \frac{dx}{x}.
\]

Tích phân:
\[
\frac{v^2}{2} = \ln|x| + C.
\]

Trở lại $v = y/x$:
\[
\frac{1}{2}\left(\frac{y}{x}\right)^2 = \ln|x| + C
\quad\Rightarrow\quad
y^2 = 2x^2 (\ln|x| + C).
\]

\paragraph{Ví dụ 2.}
Giải phương trình:
\[
(x^2 - xy)\,dx + (y^2 - xy)\,dy = 0.
\]

Kiểm tra tính đẳng cấp:
\[
P(tx,ty)=t^2(x^2-xy),\qquad
Q(tx,ty)=t^2(y^2-xy).
\]
→ đẳng cấp bậc 2.

Ta viết lại:
\[
\frac{dy}{dx} = -\frac{x^2 - xy}{y^2 - xy}.
\]

Đặt $v=\dfrac{y}{x}$:
\[
y=vx,\quad y' = v + xv'.
\]

Thay vào:
\[
v + xv' = -\frac{1 - v}{v^2 - v}.
\]

Rút gọn, tách biến rồi tích phân (bỏ qua chi tiết dài):
\[
\ln\left|\frac{v-1}{v}\right| = \ln|x| + C.
\]

Trở lại biến $y$:
\[
\frac{y-x}{y} = Cx
\quad\Rightarrow\quad
y = \frac{x}{1 - Cx}.
\]


\subsection{Phương trình tuyến tính cấp một}
Dạng tổng quát:
\[
y' + p(x)y = q(x).
\]

\subsubsection*{Công thức nghiệm}
\[
y = e^{-\int p(x)\,dx}\left( C + \int q(x)e^{\int p(x)\,dx}\,dx \right).
\]

\subsubsection*{Ví dụ}
\[
y' - y = e^{2x}.
\]
\textbf{Giải:}
\[
p(x) = -1,\quad q(x)=e^{2x}.
\]
\[
\mu(x)=e^{\int -1 dx}=e^{-x}.
\]
\[
y = e^{x}\left( C + \int e^{-x}e^{2x}dx\right)
= e^{x}\left( C + \int e^{x}dx \right)
= e^{x}(C + e^{x}).
\]

\subsection{Phương trình Bernoulli}
\[
y' + p(x)y = q(x)y^{a}, \quad a\neq 0,1.
\]
Đặt:
\[
u = y^{1-a}.
\]
Suy ra PTVP tuyến tính theo $u$.

\subsubsection*{Ví dụ}
\[
y' + 2y = y^2.
\]
Đặt $u = y^{-1}$.  
\[
u' = -y^{-2}y' = -y^{-2}(y^2 - 2y) = -1 + 2u.
\]
\[
u' - 2u = -1.
\]
Giải tuyến tính:
\[
u = Ce^{2x} + \frac{1}{2}.
\]
Suy ra:
\[
y = \frac{1}{u} = \frac{1}{Ce^{2x} + \frac{1}{2}}.
\]

\subsection{Phương trình vi phân toàn phần}

Xét phương trình:
\[
P(x,y)\,dx + Q(x,y)\,dy = 0.
\]

\subsubsection*{Điều kiện toàn phần}
Phương trình là toàn phần nếu:
\[
\frac{\partial P}{\partial y} = \frac{\partial Q}{\partial x}.
\]

Khi đó tồn tại hàm thế $u(x,y)$ sao cho:
\[
du = P\,dx + Q\,dy,
\quad\text{và nghiệm được cho bởi } u(x,y)=C.
\]

\subsubsection*{Công thức tìm hàm thế $u(x,y)$ (theo slide Giải tích 2 trang 50)}
Nếu miền $D=\mathbb{R}^2$ hoặc miền đơn liên, ta có thể xác định $u(x,y)$ theo hai công thức sau:

\paragraph{Công thức 1 (đi theo hướng ngang rồi dọc).}
\[
\boxed{
u(x,y) = \int_{x_0}^{x} P(x,y_0)\,dx
           + \int_{y_0}^{y} Q(x,y)\,dy
}
\]

\paragraph{Công thức 2 (đi theo hướng dọc rồi ngang).}
\[
\boxed{
u(x,y) = \int_{y_0}^{y} Q(x_0,y)\,dy
           + \int_{x_0}^{x} P(x,y)\,dx
}
\]

Hai công thức trên tương đương do vi phân $P\,dx + Q\,dy$ là vi phân toàn phần, nên giá trị $u(x,y)$ không phụ thuộc đường đi.

\bigskip

\subsubsection*{Cách làm thực hành}
Ngoài hai công thức trên, cách thông dụng nhất để xây dựng $u(x,y)$ là:

\paragraph{Bước 1.} Tích phân theo $x$:
\[
u(x,y) = \int P(x,y)\,dx + \varphi(y).
\]

\paragraph{Bước 2.} Lấy đạo hàm theo $y$ rồi so sánh với $Q(x,y)$:
\[
u_y = \frac{\partial}{\partial y}\left(\int P(x,y)\,dx\right) + \varphi'(y)
= Q(x,y).
\]
Từ đó suy ra $\varphi'(y)$ và tích phân để tìm $\varphi(y)$.

\bigskip

\subsubsection*{Ví dụ minh họa}
Giải phương trình:
\[
(2xy + 1)\,dx + (x^2 + \cos y)\,dy = 0.
\]

Ta có:
\[
P=2xy+1,\qquad Q=x^2+\cos y.
\]
\[
P'_y = 2x = Q'_x \quad\Rightarrow\quad \text{phương trình toàn phần.}
\]

Áp dụng công thức:
\[
u(x,y)=\int P(x,y)\,dx + \varphi(y)
      = \int (2xy+1)\,dx
      = x^2 y + x + \varphi(y).
\]

Tính đạo hàm theo $y$:
\[
u_y = x^2 + \varphi'(y).
\]

So với $Q(x,y)=x^2+\cos y$ suy ra:
\[
\varphi'(y)=\cos y \quad\Rightarrow\quad \varphi(y)=\sin y.
\]

Hàm thế:
\[
u(x,y)=x^2 y + x + \sin y.
\]

Nghiệm:
\[
x^2 y + x + \sin y = C.
\]


\section{Phương trình vi phân cấp hai giảm cấp}

\subsection{Khuyết $y'$}
Dạng:
\[
F(x, y, y'') = 0.
\]
Đặt $p = y'$, ta có:
\[
F(x, y, p') = 0.
\]

\subsection{Khuyết $y$}
\[
F(x, y', y'') = 0.
\]

\subsection{Khuyết $x$}
\[
F(y, y', y'') = 0.
\]
Đặt $p = y'$.
\[
y'' = \frac{dp}{dx} = \frac{dp}{dy}y'.
\]

\subsubsection*{Ví dụ}
Giải:
\[
y'' = yy'.
\]
Đặt $p=y'$:
\[
p\frac{dp}{dy} = yp.
\]
\[
\frac{dp}{dy}=y.
\]
\[
p = \frac{y^2}{2} + C.
\]
Suy ra:
\[
y' = \frac{y^2}{2} + C.
\]
Phương trình cấp một tách được biến.

\section{PTVP tuyến tính cấp hai - hệ số hằng}

Xét phương trình:
\[
y'' + a_1 y' + a_2 y = 0.
\]

Phương trình đặc trưng:
\[
k^2 + a_1 k + a_2 = 0.
\]

Tùy vào nghiệm của phương trình đặc trưng ta có ba dạng nghiệm.

\subsection{Hai nghiệm thực phân biệt}
Nếu phương trình đặc trưng có hai nghiệm thực khác nhau $k_1 \neq k_2$:
\[
y = C_1 e^{k_1 x} + C_2 e^{k_2 x}.
\]

\subsubsection*{Ví dụ}
Giải phương trình:
\[
y'' - 5y' + 6y = 0.
\]
\textbf{Giải:}

Phương trình đặc trưng:
\[
k^2 - 5k + 6 = 0 \Rightarrow k = 2,\; 3.
\]

Do đó nghiệm tổng quát:
\[
y = C_1 e^{2x} + C_2 e^{3x}.
\]

\bigskip

\subsection{Nghiệm kép}
Nếu phương trình đặc trưng có nghiệm kép $k$:
\[
y = (C_1 + C_2 x)e^{kx}.
\]

\subsubsection*{Ví dụ}
Giải phương trình:
\[
y'' - 4y' + 4y = 0.
\]
\textbf{Giải:}

Phương trình đặc trưng:
\[
k^2 - 4k + 4 = 0 \Rightarrow (k-2)^2 = 0.
\]

Nghiệm kép $k=2$, do đó:
\[
y = (C_1 + C_2 x)e^{2x}.
\]

\bigskip

\subsection{Nghiệm phức}
Nếu phương trình đặc trưng có hai nghiệm phức liên hợp
\[
k = \alpha \pm i\beta,
\]
thì nghiệm có dạng:
\[
y = e^{\alpha x}\left(C_1 \cos(\beta x) + C_2 \sin(\beta x)\right).
\]

\subsubsection*{Ví dụ}
Giải phương trình:
\[
y'' + 2y' + 5y = 0.
\]
\textbf{Giải:}

Phương trình đặc trưng:
\[
k^2 + 2k + 5 = 0.
\]
\[
k = -1 \pm 2i.
\]

Suy ra nghiệm:
\[
y = e^{-x}\left(C_1 \cos 2x + C_2 \sin 2x \right).
\]

\section{Phương trình vi phân tuyến tính cấp hai không thuần nhất}

Xét phương trình:
\[
y'' + a_1 y' + a_2 y = f(x).
\]

Nghiệm của phương trình tổng quát gồm hai phần:
\[
y = y_h + y_p,
\]
trong đó:
\begin{itemize}
    \item $y_h$ là nghiệm tổng quát của phương trình thuần nhất:
    \[
    y'' + a_1 y' + a_2 y = 0,
    \]
    \item $y_p$ là một nghiệm riêng của phương trình không thuần nhất.
\end{itemize}

\subsection{Phương pháp hệ số bất định}
Xét phương trình:
\[
y'' + a_1 y' + a_2 y = f(x).
\]
Nghiệm tổng quát:
\[
y = y_h + y_p.
\]

Dưới đây là các dạng quan trọng của $f(x)$ và cách đoán nghiệm riêng $y_p$, chia theo
\[
f(x)=e^{\alpha x}P_n(x), \qquad 
f(x)=e^{\alpha x}\left(P_m(x)\cos\beta x + Q_n(x)\sin\beta x\right),
\]
với đa thức $P_m,Q_n$.

\bigskip

%%%%%%%%%%%%%%%%%%%%%%%%%%%%%%%%%%%%%%%%%%%%%%%%%%%%%%%
\subsubsection{Trường hợp $f(x)=e^{\alpha x}P_n(x)$}

\paragraph{TH1 — $\alpha$ không phải nghiệm đặc trưng.}

Giải:
\[
y'' - 3y' + 2y = 4x^2 + 3x - 1.
\]

Đặc trưng:
\[
k^2 - 3k + 2 = 0 \Rightarrow k=1,2.
\]
$\alpha=0$ không trùng.

Đoán nghiệm riêng:
\[
y_p = Ax^2 + Bx + C.
\]

Thế vào → giải hệ → 
\[
A=2,\; B=\tfrac72,\; C=\tfrac32.
\]

\[
\Rightarrow y_p = 2x^2 + \frac72 x + \frac32.
\]

\bigskip

\paragraph{TH2 — $\alpha$ là nghiệm đơn của đặc trưng.}

Giải:
\[
y'' - 3y' + 2y = e^{x}(x^2 - 3x + 2).
\]

Đặc trưng:
\[
k=1,2.
\]

$\alpha=1$ trùng nghiệm đơn → nhân thêm $x$.

\[
y_p = x e^{x}(Ax^2 + Bx + C).
\]

Giải hệ thu được:
\[
A=1,\; B=-1,\; C=0.
\]

\[
y_p = x e^{x}(x^2 - x).
\]

\bigskip

\paragraph{TH3 — $\alpha$ là nghiệm kép của đặc trưng.}

Giải:
\[
y'' - 2y' + y = e^{x}(4x+1).
\]

Đặc trưng:
\[
(k-1)^2=0.
\]

Vì $\alpha=1$ trùng nghiệm kép → nhân $x^2$:

\[
y_p = x^2 e^{x}(Ax + B).
\]

Giải hệ:
\[
A=2,\quad B=-1.
\]

\[
y_p = x^2 e^{x}(2x-1).
\]

\bigskip

%%%%%%%%%%%%%%%%%%%%%%%%%%%%%%%%%%%%%%%%%%%%%%%%%%%%%%%
\subsubsection{Trường hợp $f(x)=e^{\alpha x}(P_m(x)\cos\beta x + Q_n(x)\sin\beta x)$}

\paragraph{TH4 — $\alpha+i\beta$ không phải nghiệm đặc trưng.}

Giải:
\[
y'' - y = e^{x}\big[(2x+1)\cos 3x + (x^2 - 1)\sin 3x\big].
\]

Đặc trưng:
\[
k^2 - 1 = 0 \Rightarrow k=\pm 1.
\]

\[
\alpha + i\beta = 1 + 3i \notin\{\pm 1\}.
\]

Do đó đoán nghiệm:
\[
y_p = e^{x}\big[(Ax^2+Bx+C)\cos 3x + (Dx^2+Ex+F)\sin 3x\big].
\]

(Không giải hệ do quá dài — giáo trình chỉ yêu cầu đúng dạng thử.)

\bigskip

\paragraph{TH5 — $\alpha+i\beta$ là nghiệm đơn của đặc trưng.}

Giải:
\[
y'' - 2y' + 2y = e^{x}(x^2 + 1)\cos x.
\]

Đặc trưng:
\[
k^2 - 2k + 2=0 \Rightarrow k = 1 \pm i.
\]

Ta có:
\[
\alpha=1,\; \beta=1,\; \alpha+i\beta = 1+i \text{ trùng nghiệm đặc trưng}.
\]

→ phải nhân thêm $x$:

\[
y_p = x e^{x}\big[(Ax^2+Bx+C)\cos x + (Dx^2+Ex+F)\sin x\big].
\]

(Như trên, giáo trình chỉ yêu cầu dạng đúng.)

\bigskip

\subsection{Phương pháp biến thiên hằng số (Lagrange)}

Áp dụng cho mọi hàm $f(x)$ (tổng quát hơn hệ số bất định).

\subsubsection*{Ý tưởng}
Nếu $y_1, y_2$ là hai nghiệm độc lập của phương trình thuần nhất, ta tìm nghiệm riêng dạng:
\[
y_p = C_1(x) y_1 + C_2(x) y_2.
\]

Đặt điều kiện:
\[
\begin{cases}
C_1'(x) y_1 + C_2'(x) y_2 = 0,\\
C_1'(x) y_1' + C_2'(x) y_2' = f(x).
\end{cases}
\]

Giải hệ để tìm $C_1', C_2'$, rồi lấy tích phân.

\subsubsection*{Ví dụ 3 (biến thiên hằng số)}
Giải phương trình:
\[
y'' - y = \frac{1}{1+e^{-x}}.
\]

\textbf{Giải:}

Ta có nghiệm thuần nhất:
\[
y_h = C_1 e^x + C_2 e^{-x}.
\]
\[
y_1 = e^x,\qquad y_2 = e^{-x}.
\]

Thiết lập hệ:
\[
\begin{cases}
C_1' e^x + C_2' e^{-x} = 0,\\[6pt]
C_1' e^x - C_2' e^{-x} = \dfrac{1}{1+e^{-x}}.
\end{cases}
\]

Cộng hai phương trình:
\[
2C_1' e^x = \frac{1}{1+e^{-x}}.
\]

Suy ra:
\[
C_1' = \frac{1}{2 e^x (1+e^{-x})}
      = \frac{1}{2(e^x + 1)}.
\]

Lấy tích phân:
\[
C_1(x) = -\frac12 \ln(1+e^{-x}) + K_1.
\]

Trừ hai phương trình:
\[
2C_2' e^{-x} = - \frac{1}{1+e^{-x}},
\]
\[
C_2' = -\frac{e^x}{2(1+e^x)}.
\]

Tích phân:
\[
C_2(x)= -\frac12 \ln(1+e^x) + K_2.
\]

\textbf{Nghiệm riêng:}
\[
y_p = C_1(x)e^x + C_2(x)e^{-x}
= \frac12 x.
\]

\textbf{Nghiệm tổng quát:}
\[
y = C_1 e^x + C_2 e^{-x} + \frac{x}{2}.
\]

\bigskip

\subsection{So sánh hai phương pháp}

\begin{itemize}
    \item \textbf{Hệ số bất định}: nhanh, đẹp, nhưng chỉ dùng được khi $f(x)$ thuộc các dạng “đoán được”.
    \item \textbf{Biến thiên hằng số}: giải được mọi trường hợp, nhưng tính toán dài hơn.
\end{itemize}


\end{document}